\documentclass[../../../include/open-logic-section]{subfiles}

\begin{document}

\olfileid[hi]{sth}{ordinals}{vn}	
\olsection[वॉन नोयमान की रचना]{वॉन नोयमान द्वारा क्रमसूचकों की रचना}

\olref[sth][ordinals][iso]{thm:woalwayscomparable} एक विचार सुझाता है। हम
सुव्यवस्थाओं के प्रतिनिधि के रूप में कुछ वस्तुएँ प्रस्तुत कर सकते हैं जिन्हें
\emph{क्रम-प्रकार} कहते हैं। सुव्यवस्था $\tuple{A, <}$ के क्रम-प्रकार के
लिए $\ordtype{A, <}$ लिखकर हम निम्न दो सिद्धांत पाना चाहेंगे:
\begin{align*}
	\ordtype{A, <} = \ordtype{B, \lessdot} & 
	\text{ तभी और केवल तभी जब } \ordeq{\tuple{A, <}}{\tuple{B, \lessdot}}\\
	\ordtype{A, <} < \ordtype{B, \lessdot}&
	\text{ तभी और केवल तभी जब किसी }b \in B\text{ के लिए }
	\ordeq{\tuple{A, <}}{\tuple{B_b, \lessdot_b}}
\end{align*}
इसके अतिरिक्त हम क्रम-प्रकारों को \emph{कुछ समुच्चयों के रूप में} प्रस्तुत
करने की आशा कर सकते हैं, जैसे प्राकृतिक संख्याओं को कुछ समुच्चयों के रूप
में प्रस्तुत कर सकते हैं।

ऐसा करने का सबसे सामान्य ढंग---और हमारा अपनाया दृष्टिकोण---इन क्रम-प्रकारों
को कुछ \emph{विहित} सुव्यवस्थित समुच्चयों द्वारा परिभाषित करना है। इन विहित
समुच्चयों को पहले वॉन नोयमान ने प्रस्तुत किया था:

\begin{defn}
समुच्चय $A$ \emph{संक्रमणशील} तभी और केवल तभी है जब
$(\forall x \in A)x \subseteq A$। फिर $A$ \emph{क्रमसूचक} तभी और केवल तभी
है जब $A$ संक्रमणशील हो और $\in$ द्वारा सुव्यवस्थित हो।
\end{defn}
\noindent
आगे हम क्रमसूचकों के लिए यूनानी अक्षरों का उपयोग करेंगे। परिभाषा से तुरंत
निष्पन्न होता है कि यदि $\alpha$ कोई क्रमसूचक है, तो
$\tuple{\alpha, \in_\alpha}$ सुव्यवस्था है, जहाँ
$\in_\alpha = \Setabs{\tuple{x, y} \in \alpha^2}{x \in y}$। अतः संकेतन
का थोड़ा दुरुपयोग करके हम बस कह सकते हैं कि $\alpha$ \emph{स्वयं} सुव्यवस्था
है।

हमारे पहले कुछ क्रमसूचक ये हैं:
\[
	\emptyset, \{\emptyset\}, 
	\{\emptyset, \{\emptyset\}\}, 
	\{\emptyset, \{\emptyset\}, \{\emptyset, \{\emptyset\}\}\}, \ldots
\]
आप देखेंगे कि ये वही पहले कुछ क्रमसूचक हैं जिनसे हम अनंतता के अपने अभिगृहीत,
अर्थात वॉन नोयमान की $\omega$ की परिभाषा, में मिले थे
(\olref[sth][z][infinity-again]{sec} देखें)। यह संयोग नहीं। वॉन नोयमान की
क्रमसूचकों की परिभाषा प्राकृतिक संख्याओं को क्रमसूचक मानती है, किंतु
परिमितेतर क्रमसूचकों को भी स्थान देती है।

सदा की तरह अब हम पूछ सकते हैं: क्या ये क्रमसूचक \emph{हैं}? अथवा वॉन
नोयमान ने हमें केवल कुछ ऐसे समुच्चय दिए हैं जिन्हें हम क्रमसूचक \emph{मान}
सकते हैं? इस प्रश्न पर होने वाली चर्चाएँ \olref[sfr][rel][ref]{sec},
\olref[sfr][arith][ref]{sec}, \olref[sfr][infinite][dedekindsproof]{sec}
और \olref[sth][z][nat]{sec} में हुई चर्चाओं जैसी हैं, इसलिए हम इस बात को
और नहीं खींचेंगे। आगे हम ``क्रमसूचक'' का उपयोग बस ``वॉन नोयमान क्रमसूचक''
के अर्थ में करेंगे।

\end{document}
